
\medskip
\phantomsection\label{rem:sga3-VIB-7.6.1}%
\noindent\textbf{Remark 7.6.1.---}
The equivalent conditions of
\hyperref[prop:sga3-VIB-7.6]{7.6} plainly imply that the sheaf
\(F\) under consideration is representable. The converse is
false in general:% 
\footnote{\label{note:sga3-VIB-78}Editor's note: the argument
in the original has been expanded in what follows.}
for example, if \(k\) has characteristic \(0\), let
\(G=\mathbb G_{a,k}\), and let
\[
  f:\operatorname{Spec}k\longrightarrow G
\]
be the morphism defined by the point \(1\) of \(G(k)\). Then
\(F\) is represented by the constant \(k\)-group
\(\mathbb Z_k\), whereas \(\Gamma_G(f)=G\); hence the
monomorphism
\[
  \mathbb Z_k\hookrightarrow G
\]
is not surjective.

For simplicity, suppose now that the hypotheses of the special
case in \hyperref[prop:sga3-VIB-7.6]{7.6} hold, and suppose
that \(F\) is representable. By EGA~IV\(_3\), 8.14.2, the
sheaf \(F\) is locally of finite presentation over \(k\).
The question is therefore whether the dominant monomorphism
\[
  F\longrightarrow\Gamma_G(f)
\]
is an isomorphism, or, equivalently, a closed immersion. By
\hyperref[cor:sga3-VIB-1.4.2]{1.4.2}, this is so if \(F\) is
quasi-compact; and by
\hyperref[sga3:VIA:0.5.1]{VI\(_A\), 0.5.1}, it is so if \(F\)
is connected. In particular, by
\hyperref[cor:sga3-VIB-7.2.1]{7.2.1}, it is so when \(X\) is
connected and \(f(X)\) contains the unit element of \(G\).

\medskip
\phantomsection\label{lem:sga3-VIB-7.7}%
\noindent\textbf{Lemma 7.7.---}
\emph{Let \(k\) be an algebraically closed field, let \(G\) be
a \(k\)-group locally of finite type, let \(X\) be a
geometrically reduced \(k\)-scheme locally of finite type, let
\(f:X\to G\) be a \(k\)-morphism, and let \(H\) be a group
subscheme of \(G\) such that \(H\subset f(X)\). Put}
\phantomsection\label{form:sga3-VIB-7.7-Gamma-prime}%
\[
  \Gamma'=\bigcup_{n\geq1}f^n(X^n),
  \qquad
  \Gamma'_0=\Gamma'\cap G(k),
  \qquad
  H_0=H(k),
\]
\emph{and suppose that \(H_0\) has finite index in
\(\Gamma'_0\). Then there exists an integer \(m\) such that}
\[
  f^m(X^m)=\Gamma_G(f)
\]
\emph{(cf.\ \hyperref[prop:sga3-VIB-7.6]{7.6}), and
\(\Gamma_G(f)\) is the union of finitely many translates of
\(H\).}

For every \(n\geq1\), the subset \(f^n(X^n)\) of \(G\) is
ind-constructible (EGA~IV\(_1\), 1.9.5\textup{(viii)}).
The same is therefore true of \(\Gamma'\). Since \(G\) is a
Jacobson scheme, \(\Gamma'_0\) is dense in \(\Gamma'\). By
hypothesis, there are finitely many points
\[
  a_1,\ldots,a_r\in\Gamma'_0
\]
such that
\[
  \Gamma'_0=a_1H_0\cup\cdots\cup a_rH_0.
\]
Consequently,
\phantomsection\label{form:sga3-VIB-7.7-closure-chain}%
\[
\begin{aligned}
  \Gamma_G(f)
    &=\overline{\Gamma'}
     =\overline{\Gamma'_0}
     =\overline{a_1H_0\cup\cdots\cup a_rH_0}\\
    &=\overline{a_1H_0}\cup\cdots\cup\overline{a_rH_0}\\
    &=a_1\overline{H_0}\cup\cdots\cup
      a_r\overline{H_0},
\end{aligned}
\]
where the last equality follows because translation by each
\(a_i\) is a homeomorphism of \(G\) onto itself. Thus
\[
  \Gamma_G(f)=a_1H\cup\cdots\cup a_rH.
\]

It is also clear that there exists an integer \(p\) such that
every \(a_i\), for \(1\leq i\leq r\), belongs to
\(f^p(X^p)\). Finally, since \(H\subset f(X)\), one has
\[
  a_iH\subset f^{p+1}(X^{p+1})
\]
for every \(i\). Hence
\[
  f^{p+1}(X^{p+1})=\Gamma_G(f).
\]
